% Build with: pdflatex "Chris Cameron - Anthology resume.tex"
\documentclass[11pt,a4paper]{article}

\usepackage[T1]{fontenc}
\usepackage[utf8]{inputenc}
\usepackage{lmodern}
\usepackage[margin=2cm]{geometry}
\usepackage{enumitem}
\usepackage{titlesec}
\usepackage{xcolor}
\usepackage[hidelinks]{hyperref}
\usepackage{microtype}
\hypersetup{pdftitle={Chris Cameron - Anthology resume},pdfauthor={Chris Cameron}}

\definecolor{accent}{RGB}{30,60,110}
\definecolor{muted}{RGB}{85,85,85}
\pagestyle{empty}
\setlength{\parindent}{0pt}
\setlength{\parskip}{6pt}
\linespread{1.10}
\hyphenpenalty=10000
\emergencystretch=2em
\raggedright

\titleformat{\section}{\large\bfseries\color{accent}}{}{0pt}{}[\vspace{3pt}\titlerule]
\titlespacing*{\section}{0pt}{16pt}{8pt}
\setlist[itemize]{leftmargin=1.2em,itemsep=5pt,topsep=4pt,parsep=0pt}

\begin{document}

{\fontsize{28}{32}\selectfont\bfseries Chris Cameron\par}
\vspace{2pt}
\href{mailto:chrisjamescameron1@gmail.com}{chrisjamescameron1@gmail.com}
\enspace$\cdot$\enspace
\href{https://github.com/ccam80}{github.com/ccam80}

\vspace{8pt}
AI-native researcher, developer, and mechatronics engineer. PhD researcher
in electrophysiology (completion expected early 2027) and creator of Cubie.
Part-time lecturer in mechatronics engineering, including computer architecture, C programming,
circuit design, and embedded systems.

\section{Cubie --- GPU numerical integration}
{\small\color{muted}\href{https://github.com/cubiepy/cubie}{github.com/cubiepy/cubie}\par}
I built Cubie, an open-source Python ODE solver for consumer NVIDIA GPUs, to support
electrophysiology simulations for a real-time application. Current benchmarks
(draft, unpublished) show \textbf{10,000$\times$ speedup}
over CPU implementations (SciPy, MATLAB) and \textbf{100$\times$}
over the nearest GPU competitor, DiffEqGPU.jl, on real-world problems.

Once CUDA-shaped algorithms were in place and tiling and memory placement
were addressed, further speedups were achieved by:

\begin{itemize}
  \item using NVIDIA profiling tools to trace bottlenecks and stalls
    to individual instructions;
  \item rolling/unrolling loops based on measured instruction counts against the instruction cache;
  \item decreasing occupancy to reduce L2 cache misses (this one was surprising!);
  \item ordering and fusing operations at a whole-kernel level to reduce live ranges and reduce spilling; and
  \item placing buffers heuristically to trade off guaranteed-L1 shared memory
    against possible-L2/VRAM local memory based on working set size.
\end{itemize}

Cubie goes fast, but it is a scientific tool. The codebase is complex and correctness is critical.
To maintain correctness, I had to:

\begin{itemize}
  \item test numerical results against known references and custom CPU harnesses;
  \item review every pull request line-by-line (AI agents are prolific but often incorrect); and
  \item build approval processes including a ratcheting benchmark per PR and automated correctness checks.
\end{itemize}

\section{Open-source contributions}
\textbf{numba-cuda / numba-cuda-mlir}\par
Many merged and open pull requests address compiler performance, typing and
lowering correctness, AST transforms, and kernel performance controls.

\textbf{SymPy and lcapy} --- minor edits.

\end{document}
